\subsection{Fiber bundles associated with a principal fibration}

6.5.1. Let \(\lambda = (\mathbf{P}, \mathbf{G}, \mathbf{B}, \pi)\) be a principal fibration. Let F be a variety on which the group G acts on the left; we denote \((g, y) \mapsto g \cdot y\) the operation law of G on F. The group G acts on the right on \(\mathbf{P} \times \mathbf{F}\) by the formula \((x, f) \cdot g = (x \cdot g, g^{-1} \cdot f)\) ; the equivalence relation defined by G in \(\mathbf{P} \times \mathbf{F}\) is regular; the quotient \(\mathbf{P} \times^{\mathbf{G}}\mathbf{F} = (\mathbf{P} \times \mathbf{F}) / \mathbf{G}\) is endowed with a manifold structure.

Let E be a manifold. We say that E is equipped with a structure of fiber bundle associated with λ with typical fiber F when we are given a morphism ρ : P × F → E satisfying the following property:

(As) We have \(\rho(x \cdot g,g^{-1} \cdot f)=\rho(x,f)\) for \(x\in\mathbf{P}\) , \(f\in\mathbf{F}\) , \(g\in\mathbf{G}\) , and the map \(\bar{\rho}:\mathbf{P}\times^{\mathbf{G}}\mathbf{F}\to\mathbf{E}\) deduced from \(\rho\) by passing to the quotient is an isomorphism of varieties.

It amounts to the same thing to say that \((\mathbf{P} \times \mathbf{F}, \mathbf{G}, \mathbf{E}, \rho)\) is a principal fibration. The map \(\rho\) (or sometimes the map \(\bar{\rho}\) ) is called the frame map of E, and is denoted \((x, f) \mapsto x \cdot f\) ; we have

\[
(x. g). f = x. (g. f), \quad \text {   for   } x \in \mathbf {P}, g \in \mathbf {G} \text {   and   } f \in \mathbf {F}.
\]

The data of \(\lambda\) and F determines E up to a unique isomorphism; in particular, one may take for E the variety P × \(^{\text{G}}\) F itself; it is said to be the associated fiber bundle to \(\lambda\) of fiber type F, and it is denoted \(\lambda(\text{F})\) .

6.5.2. Let E be a fiber bundle associated with \(\lambda\) and with typical fiber F. There exists a unique morphism \(\pi_{\mathbf{E}}\) from E into B such that \(\pi_{\mathbf{E}}(x,f)=\pi(x)\) if \(x\in\mathbf{P},f\in\mathbf{F}\) ; the triple (E, B, \(\pi_{\mathbf{E}}\) ) is a fibration; if B and F are separated, then E is separated.

Let \(b \in \mathbf{B}\) , and let \(\mathbf{F}_b = \pi_{\mathbf{E}}^{-1}(b)\) ; it is a closed submanifold of E. If \(x \in \mathbf{P}\) is such that \(\pi(x) = b\) , let \(\theta_x: \mathbf{F} \to \mathbf{F}_b\) be the map defined by \(\theta_x(f) = x.f\) ; it is an isomorphism of manifolds. Moreover, for every \(g \in \mathbf{G}\) , one has \(\theta_{x.g} = \theta_x \circ \rho_g\) , where \(\rho_g\) denotes the automorphism \(f \mapsto g.f\) of F. Suppose

F is equipped with a structure s of arbitrary species \(\Sigma\) (cf.~Ens., chap.~IV, §1, n° 4) and suppose that s is invariant under G; there then exists on \(\mathbf{F}_{b}\) a structure \(s_{b}\) of species \(\Sigma\) and exactly one such that the \(\theta_{x}: F \to F_{b}\) are isomorphisms; one obtains it by transporting s by means of one of the \(\theta_{x}\) (loc. cit., no. 5).

If \(s\) is a section of \(\mathbf{P}\) above an open set \(\mathbf{U}\) of \(\mathbf{B}\) , the map \((b, f) \mapsto s(b)\) . \(f\) is an isomorphism of \(\mathbf{U} \times \mathbf{F}\) onto \(\pi_{\mathbf{E}}^{-1}(\mathbf{U})\) .

6.5.3. Examples:

\begin{enumerate}
\def\labelenumi{(\alph{enumi})}
\tightlist
\item
  Let \(\lambda = (\mathbf{B} \times \mathbf{G}, \mathbf{B}, \mathrm{pr}_1)\) , let \(\mathbf{E} = \mathbf{B} \times \mathbf{F}\) , and let
\end{enumerate}

\[
\rho : (\mathbf {B} \times \mathbf {G}) \times \mathbf {F} \rightarrow \mathbf {E}
\]

be the map \((b, g, f) \mapsto (b, g \cdot f)\) . We thus obtain on \(\mathbf{B} \times \mathbf{F}\) a fiber bundle structure associated with \(\lambda\) with typical fiber \(\mathbf{F}\) , which is said to be trivial.

\begin{enumerate}
\def\labelenumi{(\alph{enumi})}
\setcounter{enumi}{1}
\item
  Let \(\lambda = (\mathbf{P},\mathbf{G},\mathbf{B},\pi)\) and \(\lambda' = (\mathbf{P}',\mathbf{G}',\mathbf{B}',\pi')\) be two principal fibrations. Let F (resp. F') be a variety on which G (resp. G') acts on the left; and let E (resp. E') be a fiber space associated with \(\lambda\) (respectively to \(\lambda'\) ) with fiber type F (respectively F'). The group \(\mathbf{G} \times \mathbf{G}'\) acts on \(\mathbf{F} \times \mathbf{F}'\) by \((g, g') \cdot (f, f') = (g \cdot f, g' \cdot f')\) . The map \((\mathbf{P} \times \mathbf{P}') \times (\mathbf{F} \times \mathbf{F}') \to \mathbf{E} \times \mathbf{E}'\) is the product of the frame maps \(\mathbf{P} \times \mathbf{F} \to \mathbf{E}\) and \(\mathbf{P}' \times \mathbf{F}' \to \mathbf{E}'\) and endows \(\mathbf{E} \times \mathbf{E}'\) with a structure of fiber bundle associated with \(\lambda \times \lambda'\) with typical fiber \(\mathbf{F} \times \mathbf{F}'\) .
\item
  With the notations of (b), if one supposes that \(B' = B\) , one defines in the same way on \(E \times_{B} E'\) a fiber bundle structure associated with \(\lambda \times_{B} \lambda'\) with typical fiber \(F \times F'\) .
\end{enumerate}

6.5.4. The notations being those of no. 6.5.1, let \(h: \mathbf{B}' \to \mathbf{B}\) be a morphism, let \(\lambda' = \mathbf{B}' \times_{\mathbf{B}} \lambda\) , and let \(\mathbf{E}' = \mathbf{B}' \times_{\mathbf{B}} \mathbf{E}\) . If \(\mathbf{P}' = \mathbf{B}' \times_{\mathbf{B}} \mathbf{P}\) , define a map \(\mathbf{P}' \times \mathbf{F} \to \mathbf{E}'\) by setting \((b', x).f = (b', x.f)\) ; one thus endows \(\mathbf{E}'\) with a structure of fiber bundle associated with \(\lambda'\) with fiber type F; it is called the inverse image by \(h\) of the structure given on E.

6.5.5. Let \(\lambda = (\mathbf{P}, \mathbf{G}, \mathbf{B}, \pi)\) be a principal fibration, and let E (resp. E') be a fiber space associated with \(\lambda\) with fiber type F (resp. F'). Let \(u: \mathbf{F} \to \mathbf{F}'\) be a morphism compatible with the operations of G (that is, such that \(u(g.f) = g \cdot u(f)\) for \(f \in \mathbf{F}\) , \(g \in \mathbf{G}\) ). There then exists a unique morphism \(\bar{u}: \mathbf{E} \to \mathbf{E}'\) such that one has \(\bar{u}(x.f) = x \cdot u(f)\) for \(x \in \mathbf{P}\) , \(f \in \mathbf{F}\) . If \(u\) is an immersion (resp. a submersion, a subimmersion), the same is true of \(\bar{u}\) .

In particular, suppose that a group variety H operates on the right on F in such a way that \(g.(f.h) = (g.f).h\) for \(g \in G\) , \(f \in F\) , \(h \in H\) . Every \(h \in H\) then defines an automorphism \(u_h\) of F compatible with the operations of G, whence an automorphism \(\bar{u}_h\) of E; the group H operates on the right on E by \((y,h) \mapsto \bar{u}_h(y)\) .

6.5.6. Let \(\lambda = (\mathbf{P},\mathbf{G},\mathbf{B},\pi)\) be a principal fibration, and let E be a fiber space associated with \(\lambda\) with fiber type F. Let \(s: \mathbf{B} \to \mathbf{E}\) be a section of E. For every \(x \in \mathbf{P}\) , there exists a unique element \(\sigma(x) \in \mathbf{F}\) such that \(s(\pi(x)) = x \cdot \sigma(x)\) . The map \(\sigma: \mathbf{P} \to \mathbf{F}\) thus defined is a morphism of varieties, and satisfies the identity:

\[
\sigma (x. g) = g ^ {- 1}. \sigma (x).\tag{*}
\]

The map \(s \mapsto \sigma\) is a bijection of the set of sections of E onto the set of morphisms of P into F which satisfy the identity (*).

6.5.7. Let \(\lambda = (\mathbf{P}, \mathbf{G}, \mathbf{B}, \pi)\) and \(\lambda' = (\mathbf{P}', \mathbf{G}, \mathbf{B}, \pi')\) be two principal fibrations with the same base \(\mathbf{B}\) and of the same structural group \(\mathbf{G}\) . The principal fibration \(\lambda \times_{\mathbf{B}} \lambda' = (\mathbf{P} \times_{\mathbf{B}} \mathbf{P}', \mathbf{G} \times \mathbf{G}, \mathbf{B}, (\pi, \pi')_{\mathbf{B}})\) has as its structure group \(\mathbf{G} \times \mathbf{G}\) . Let us make \(\mathbf{G} \times \mathbf{G}\) operate on the left on \(\mathbf{G}\) by the formula:

\[
(g, g ^ {\prime}) \cdot g _ {1} = g \cdot g _ {1} \cdot g ^ {- 1},
\]

and let E be the fiber bundle associated with \(\lambda \times_{\mathbf{B}} \lambda'\) with fiber type G (endowed with the operation law defined above). Then the sections of E correspond bijectively to the isomorphisms from P to P'. Precisely, if s is a section of E corresponding (cf.~no. 6.5.6) to the morphism \(\sigma: \mathbf{P} \times_{\mathbf{B}} \mathbf{P}' \to \mathbf{G}\) , there exists a G-B-isomorphism \(f_s: \mathbf{P} \to \mathbf{P}'\) such that \(\sigma(x, f_s(x)) = e\) for all \(x \in \mathbf{P}\) ; the map \(s \mapsto f_s\) is a bijection from the set of sections of E onto the set of G-B-isomorphisms from P to P'.

